% QNLPAI 2026 — STUDENT POSTER SESSION ABSTRACT (author-written draft; <= 2 pages incl. refs, LNCS).
% DOUBLE-BLIND: no author identity; NO Acknowledgments section (CFP omits acknowledgments in the blind
% paper). AI-assistance disclosure deferred to camera-ready / OpenReview field.
% NOTE: body text is the AUTHOR's own wording; this pass only fixes spelling/caps/math formatting,
% adds \cite links + the figure. Content unchanged. Length WILL exceed 2pp -- compress next.
\documentclass[runningheads]{llncs}
\usepackage[T1]{fontenc}
\usepackage{lmodern}
\usepackage{microtype}
\usepackage{graphicx}
\usepackage{amsmath,amssymb}
\usepackage[hidelinks]{hyperref}

\begin{document}

\title{Language-Conditioned QAOA for Portfolio Selection: A Bias-Controlled Study Reveals a
Training-Horizon Confound}
\titlerunning{Language-Conditioned QAOA: A Training-Horizon Confound}
\author{Anonymous Author(s)}
\authorrunning{Anonymous}
\institute{Affiliation withheld for double-blind review}
\maketitle

\begin{abstract}
We flag a possible oversight in using large language models for portfolio optimization: we
build a language-conditioned QAOA in which an LLM reads RavenPack news into per-asset convictions that
tilt the cost Hamiltonian of a cardinality-constrained Quantum Approximate Optimization Algorithm (QAOA)
portfolio optimizer. We ask, with bias controls, whether language
adds out-of-sample value across two LLMs: qwen2.5-14B (Sept.\ 2024, released during the
evaluation window) and an out-of-sample Llama~2 (cutoff Sept.\ 2022) on the same headlines. The qwen model
beats matched random tilts by $+0.76\%$/quarter and Llama~2 $-0.23\%$/quarter ($p=0.94$), which is
indistinguishable from random. The benefit is noticed only for the model released into the test period,
consistent with training-horizon hindsight rather than prediction---a misleading alpha. The pre-cutoff
model is also smaller, a caveat we note. We contribute the pipeline, a reusable bias-control template, and
a training-horizon look-ahead warning for LLM-for-finance evaluation.
\keywords{QAOA \and Quantum optimization \and LLM-conditioning \and Look-ahead bias \and Quantum finance.}
\end{abstract}

\section{Pipeline and Method}
We focus on the use case of portfolio optimization in financial institutions. Chen et
al.~\cite{chen2023qbl} use Black--Litterman with quantum numeric views, and Lee
et al.~\cite{lee2025llmbl} with LLM views optimized classically; we question their combination of quantum
and LLM in one pipeline. We use
a fixed set of Dow~30 across 2022--24, split into 378-day train and 63-day (next-quarter) test windows at
each of 4 decision dates (2023-12-01 through 2024-09-04). We draw 400 random 14-of-30 subsets per date
(without replacement), 1600 in total. Each asset is one qubit; we cap subsets at 14
for a tractable exact statevector simulation ($2^{14}$ amplitudes), with the exact optimum over
$\binom{14}{5}=2002$ portfolios.

Each subset is reduced to 5 held assets under the Markowitz mean--variance objective (risk aversion
$q=2$; $\mu$ and $\Sigma$ from CRSP total returns with shrinkage covariance). Holding exactly $B=5$ of
$n=14$ makes $x$ binary---a combinatorial choice over the 2002 feasible portfolios, and thus a candidate
for a quantum optimizer.

The LLMs qwen2.5-14B (released Sept.\ 2024, its training plausibly spans our 2024 test window) and Llama~2
(pretraining cutoff Sept.\ 2022, before the test data window) are provided with RavenPack headlines
(tagged, timestamped) and each decide whether to nudge $\mu$; for each decision date the input is only
headlines strictly before it, so there is no look-ahead in the news itself. The output from the LLM is a conviction $v_i \in [-1,1]$. We
inject the view Black--Litterman-style into expected returns, $\mu \to \mu + s\,v$ with $s = 0.10$ (a
$\pm10\%$ annual view), not tuned to force the optimizer to change its picks. As a control, the placebo matches the tilt's size and sparsity
but randomizes which stocks (300 draws form a null), testing whether the language matters or whether any
nudge of that size would do the same.

We encode the binary ``pick 5 of 14'' as a QUBO (mean--variance score plus a ``hold exactly 5'' penalty),
mapping to an Ising model whose minimum energy is the best portfolio. We solve with QAOA at depth $p=4$
($2p=8$ variational parameters). Our objective is not quantum advantage but whether the LLM's view changed
the best portfolio at the exact optimum; the economic effect is read there, with QAOA's sampling quality
reported separately.

\section{Results}
On the same headlines, qwen2.5-14B's changed-optimum tilts beat matched random tilts by $+0.76\%$/quarter,
whereas Llama~2's do $-0.23\%$/quarter. The edge appears only for the model whose training overlaps the test window;
swapping to the out-of-sample model (market and factors held fixed) erases it---our core finding: the edge
is the training horizon, not the market. The $+0.76\%$ is the language-specific premium: qwen's
changed-instance effect ($+1.20\%$) minus matched random tilts ($+0.43\%$). The language input changes the
exact optimum in $47$--$68\%$ of instances, having an effect on selection.

Qwen's instance-level $p\approx0$ (outside the 300-draw null), positive in 3 of 4 quarters; but the 1600
instances are subsets within only 4 decision dates, so this $p$ is pseudo-replicated, not a time-series
test. Llama~2 gives $p=0.94$ (2 of 4 quarters), and the per-date effect is regime-dependent. QAOA at
depth $p=4$ recovers the true optimum only $7.7\%$ of the time, a weak exact optimizer at this depth.

\paragraph{Limitations.} Qwen (14B) is larger than Llama~2 (7B), and qwen has no published cutoff, so
contamination is inferred; the study is a small-$n$ exact simulation without transaction costs. These
restrict causal attribution but do not overturn the pattern---present for the recent model, absent for the
pre-cutoff one.

\bibliographystyle{splncs04}
\bibliography{references}

% CAMERA-READY: restore author block + Acknowledgments (advisor thanks + AI-assistance disclosure).
\end{document}
